Are decision models worth it, or can small task-specific classifiers do the same job? In our experiments the answer depends
mainly on how many labelled examples are available and, when many are available, on how well the small model is trained.

With 5 or 20 labelled examples per class, decision models used zero-shot with label names only were far ahead of every compact
classifier on all datasets. With the full training sets (70--190 examples per class) and our frozen recipes, the gap was 2.4 points on
Banking77, 5.5 points on Italian MASSIVE and not resolved on the synthetic Ottavo task. On Banking77 the recipe of the earlier exploratory
study raised Qwen3-0.6B from 0.910 to 0.936: the exploratory score was reproduced, but the claim that the small model wins was not, since
the difference from Mammoth-1 was no longer resolved. Among the compact methods evaluated, SetFit was the best with five examples per
class; decoders with a fresh classification head collapsed in that regime with both recipes tested, and needed at least 20 examples per
class.

Three further results qualify the picture. First, the effect of category descriptions varied with the model and the dataset: they helped
OpenAI Decisions on Banking77 and all decision models on the synthetic task, and elsewhere the observed effects were small and mixed. On the
public datasets our descriptions were written from labelled examples, so they are not a free resource. Second, Elephant-1 was much closer
to its own base checkpoint, Qwen3.5-9B, than to an older generalist; the remaining gap was small and larger on lexically adversarial
messages. Third, the Base Labs models were insensitive to option order and well calibrated, which matters for routing uncertain cases to a
human; OpenAI's coarse probabilities were less useful for that purpose.

For a practitioner this suggests a simple rule of thumb, valid for intent-style tasks like those tested: without a substantial labelled set,
a decision model (or a strong recent generalist) is the better tool; with thousands of labels, a well-trained small classifier can come close
to it, or show no resolved difference, and is cheaper to run at volume, but its quality depends on the training recipe, and the cost of
labelling and of maintenance has to be weighed. Whether these conclusions hold on real customer traffic is the main open question; a second
version of this report with human-annotated real messages would answer it.
